\section{Coefficient scope, norm lines, and stable lattices}
\label{sec:coefficient-scope}

We keep the coefficient and frame conventions of
Proposition~\ref{prop:frame-torsor}.  In particular, a cochain with values
in a localized coefficient ring has a uniform pole bound on its compact
source.  A cochain in a union of frame levels or root extensions is
represented at one finite level.  All vector spaces in this section are
over the finite field $k$.

\subsection{The simultaneous induction}

We specify the auxiliary coefficients needed on boundary strata.  This
also records the precise scope of the compact finiteness assertion.

\begin{definition}\label{def:compact:admissible}
An admissible compact coefficient on $A$ or $C_1$ is a finite free module
with continuous $P_n$-action satisfying the following condition on every
successive labeled height stratum.  After restriction to its exact-height
open, and after a finite joint frame change and a fixed finite Frobenius
root change, the module has a finite $P_n$-stable filtration whose graded pieces
admit $P_n$-equivariant identifications with finite direct sums of
the coefficient ring, with $P_n$ acting only on that ring.
We call these identifications constant frames; any commuting
auxiliary frame actions are retained.  The filtration splits as
a filtration of underlying modules.
The same condition is required after restriction to each dead-label
boundary disc and power-identification with its starting-height disc.
For a module on $A$ the condition is checked after pullback to the labeled
strata.

We also allow the constant frames to require a finite level of the
translation splitting torsor.  Equivalently, before that last ascent the
graded pieces may be associated to finite representations of a finite
translation quotient.  The condition must still hold on every successive
boundary stratum.  All actions, filtrations, and splittings use bounded
poles on each exact-height open.
\end{definition}

The boundary requirement is recursive: at a pair $(n,t)$ it also
tests the restrictions to the labelled discs with starting heights
$t+1,\ldots,n$.  A power identification of such a disc transports
its coefficient and group action together.  A finite translation
splitting can instead be kept as an ascent with its representation
data, for use in associated-bundle descent.

Here are the coefficients to which we will apply the definition.
The invariant line, the Lie bundles of the universal group and its
Cartier dual, their tensor and character constructions, and finite
torsion or Frobenius-divided torsion
coordinate modules have constant frames after marking the connected
group and the finite \'{e}tale quotient and splitting the relevant
finite lifting torsor.  A fixed finite construction needs only
finitely many marking coefficients.  Its linear dual has the dual
frame.  Thus these constructions and their finite tensor operations
satisfy the condition on each successive stratum.

Before translation splitting, the same coefficients can be treated
through their finite translation representations.  At connected
height greater than one the distribution algebra is local, so those
representations have filtrations by trivial lines.  At multiplicative
height the simple representations are characters; an associated
character is a summand of the regular representation of its finite
translation quotient.  The coordinate module of a split group itself
has constant frames.

Powers of conormal and canonical lines are also included.  The volume
construction trivializes the canonical line on an exact-height
stratum, and adjunction along a label hyperplane introduces its
conormal line.  The higher-height identification in
Corollary~\ref{rings:boundary-labels} transports these constructions
to the corresponding boundary disc.  The same verifications commute
with finite tensor operations, linear duals, and power transport.

For each height $s$, let
\[
 \nu_s:P_s\longrightarrow\Fp^\times\subset k^\times,
 \qquad
 \nu_s(g)=\overline{\operatorname{Nrd}(g)},
\]
be the reduced norm followed by reduction modulo $p$.  Write
$k(\nu_s^j)$ for the corresponding one-dimensional constant
$P_s$-representation.  These particular constant characters have a
geometric realization on the entire deformation disc.

For rectangular frame levels write
\[
 B_{K,V}=B_{K\times V},\qquad
 B_{\mathrm{conn},V}=\colim_{K\subset P_t}B_{K,V},\qquad
 B_{\mathrm{full}}=\colim_{K,V}B_{K,V},
\]
where $K$ and $V$ run through compact open subgroups of $P_t$ and
$\GL_m(\Zp)$, respectively.  These are algebraic unions with the cochain
convention above.

\Needspace{9\baselineskip}
\begin{theorem}[Coefficient finiteness]\label{thm:coefficient-finiteness}
Suppose $1\leq t<n<p-1$.  The following assertions hold.
\begin{enumerate}
\item Every admissible compact coefficient on $A$ or $C_1$ has finite
total continuous $P_n$-cohomology.  This includes the stable finite free
lattices constructed below.
\item For every finite joint frame level $U$, the bounded-pole groups
$H^*(P_n,B_U)$ have finite total dimension.  The assertion also
holds with a fixed power of the invariant line.  Finite direct
sums and finite-dimensional trivial coefficient factors are
allowed.
\item For every fixed compact open $V\subset\GL_m(\Zp)$,
$H^*(P_n,B_{\mathrm{conn},V})$ has finite total dimension.  There is an
integer $h$, independent of the rectangular levels used in the union,
such that $1+p^h\Zp$ acting by scalar matrices acts identically on
$H^*(P_n,B_{\mathrm{full}})$.  These conclusions persist after taking the
algebraic union of Frobenius roots, with the actions transported by
powering.

Moreover, put
$B_{\mathrm{full}}^{\mathrm{perf}}
 =\colim_r B_{\mathrm{full}}^{1/p^r}$, the algebraic root union
with common finite stages on compact cochain sources.
For every compact open $V\subset\GL_m(\Zp)$ and every $a,i$,
the finite-dimensional smooth $P_t$-representation
\[
 Q_{V}^{a,i}=
 H^a\!\left(V,H^i(P_n,B_{\mathrm{full}}^{\mathrm{perf}})\right)
\]
has a finite $P_t$-stable filtration whose successive quotients
are powers of the reduced-norm character $\nu_t$.
\end{enumerate}
All assertions use common finite stages on profinite cochain sources.
\end{theorem}

The theorem fixes the target of a distributed induction, completed in
Sections~\ref{sec:compact}--\ref{full:section}. We label a case by its
height pair $(n,t)$ and order the cases by the lexicographic measure
\begin{equation}\label{eq:compact:induction-order}
                         (n,n-t).
\end{equation}
At each height pair, the three assertions are proved in the displayed order.
Their induction dependencies and proof locations are as follows.
\begin{description}[style=nextline,leftmargin=0pt,labelsep=0pt,itemsep=4pt]
\item[Compact coefficients (1).]
Proposition~\ref{prop:compact-finiteness} uses compact finiteness on the
dead-label discs at the higher starting-height pairs
$(n,t+1),(n,t+2),\ldots$. The initial case $t=n$ is finite-coefficient
cohomology on the zero-dimensional disc.
\item[Bounded poles (2).]
Proposition~\ref{prop:pole-removal} uses the compact assertion at the
current pair, the full-frame assertions at $(n,s)$ for $s>t$, and compact
coefficients at the smaller total-height pairs $(s,t)$.
\item[Full frame (3).]
Propositions~\ref{full:joint-completion} and
\ref{full:perfected-row-characters} use the first two assertions at the
current pair and the support comparison. They include the specified
perfected individual-row filtration. Thus the row filtration at $(n,s)$
is available before the bounded-pole proof at $(n,t)$ whenever $s>t$.
\end{description}
Every appeal to an earlier case decreases \eqref{eq:compact:induction-order}.
In particular, the full-frame assertion at the current pair is never used
in its compact or bounded-pole proof. The cited propositions complete the
proof of Theorem~\ref{thm:coefficient-finiteness}.

\subsection{Norm-character coefficients}

We begin with the coefficients that occur in the smaller-disc
step of this induction. The next lemma realizes the required
constant norm characters on the universal deformation disc,
including every boundary restriction used below.

\begin{lemma}[The norm line on the deformation disc]
\label{lem:compact:norm-line}
Let $\mathcal G$ be the algebraic height-$s$, dimension-one
Barsotti--Tate group obtained from the universal formal deformation,
and let $\mathcal R_s$ be its characteristic-$p$ complete local
deformation ring.  There is a functorial height-one, dimension-one
Barsotti--Tate group
\[
                 \mathcal H_{\det}=\bigwedge^s\mathcal G.
\]
Its marked deformation is canonically constant.  The covariant action
of $P_s$ on its special group, and hence on this constant marked
deformation, is multiplication by the actual unit
$\operatorname{Nrd}(g)\in\Zp^\times$.

With the convention that actions on functions use inverse
substitution, put
\[
 \mathcal N_s=
       \Hom_{\mathcal R_s}(\omega_{\mathcal H_{\det}},\mathcal R_s).
\]
Then there is a $P_s$-equivariant identification
\[
             \mathcal N_s\simeq
                 \mathcal R_s\otimes_k k(\nu_s).
\]
Every integral tensor power of this line is admissible on every
starting-height quotient of the deformation disc and on its labeled
disc.  The assertion is preserved by the finite root transports and
the successive boundary restrictions used in
Definition~\ref{def:compact:admissible}.
\end{lemma}

The character in this identification need not be trivial on $P_s$.
Admissibility instead asks for a frame after a finite equivariant
base change: the frame can transform through the marking coordinates.
The proof supplies this frame from the first truncated determinant
group, without trivializing the original constant character.

\begin{proof}
The prime in this manuscript is odd.  The exterior-power theorem for
a height-$s$ $p$-divisible group of dimension at most one applies to
the algebraic Barsotti--Tate group, including its finite truncated
levels.  Its exterior powers have heights $\binom{s}{r}$ and, on
the dimension-one locus, dimensions $\binom{s-1}{r-1}$; the
construction and its universal alternating map commute with base
change.  Apply this for $r=s$ to obtain $\mathcal H_{\det}$ of
height and dimension one
\cite[Theorem~3.25]{HedayatzadehExterior}.  The same functorial
construction identifies its first truncated level with the top
exterior construction on $\mathcal G[p]$
\cite[Lemmas~3.23--3.24 and Theorem~3.25]{HedayatzadehExterior}.
We use these truncated
group schemes before any localization of their coefficient rings.

Here is an explicit justification of the constancy assertion.
Let $H_0$ be the special fibre of $\mathcal H_{\det}$, with the
marking induced from the fixed Honda marking of $\mathcal G$.
The Cartier dual of a height-one, dimension-one group is an
\'{e}tale group of height one.  The complete local ring
$\mathcal R_s$ is henselian
\cite[Lemma~10.153.9, Tag~04GM]{StacksProject}.
For every $r$, its finite group
$\mathcal H_{\det}^{\vee}[p^r]$ is the unique finite \'{e}tale
lift over $\mathcal R_s$ of
$H_0^{\vee}[p^r]$, with its specified special-fibre identification.
The constant lift along $k\to\mathcal R_s$ is one such lift.
The equivalence of finite \'{e}tale categories under reduction
\cite[Lemma~10.153.7, Tag~04GK]{StacksProject}
identifies the two
lifts uniquely, including their group laws, transition maps, and
homomorphisms.  These identifications are compatible for all $r$.
Dualizing gives the canonical marked identification
\[
          \mathcal H_{\det}\simeq
                   H_0\times_{\Spec k}\Spec\mathcal R_s.
\]
One can equivalently obtain this uniqueness by height-one deformation
theory: its Hodge filtration has a unique lift and its marked
morphisms are rigid; see \cite[Theorem~48, Proposition~40,
and Corollary~95]{ZinkDisplays}.  The finite \'{e}tale argument also
shows directly that all the chosen identifications are compatible
with nilpotent thickenings.  Their pullbacks give the required
identifications on every quotient and finite-level base used below.

We identify the action, rather than only its closed-point scalar.
On the rational Honda Dieudonne realization the action of the
endomorphism division algebra becomes its defining $s$-dimensional
matrix action after passage to a splitting field.  Functoriality of
the exterior construction makes the action on the top exterior
line its determinant.  By the defining identity for the reduced
norm this determinant is $\operatorname{Nrd}(g)$, and the equality
descends from the splitting field.  The normalized Frobenius in the
exterior construction changes the Frobenius of the line, not the
determinant action of an automorphism on its underlying line.
For $g\in P_s$ the determinant is an integral unit, so on the
height-one group $H_0$ it is the automorphism
$[\operatorname{Nrd}(g)]$.  Rigidity of the marked height-one lift
then identifies the semilinear transport on
$\mathcal H_{\det}$ with this same constant automorphism over
the whole parameter ring.  In particular this is not an
identification inferred merely by reducing a varying invariant-line
cocycle at the closed point.

Under inverse substitution the action on an invariant differential
of $H_0$ is multiplication by
$\overline{\operatorname{Nrd}(g)}^{-1}$.
Its dual line therefore has character $\nu_s$, proving the
displayed equivariant identification of $\mathcal N_s$.

It remains to check the finite-stage admissibility condition.
On any exact-height stratum, choose finite connected and \'{e}tale
markings sufficient for the first truncated group, and a finite
translation splitting of its connected--\'{e}tale extension.
These are exactly the permitted ascents in
Definition~\ref{def:compact:admissible}.  They give an equivariant
constant frame for $\mathcal G[p]$ on that stratum.  Functoriality
and base change for the truncated exterior construction then give
an equivariant constant frame for $\mathcal H_{\det}[p]$.
In characteristic $p$, the invariant differential module of a
$p$-divisible group is the invariant differential module of its
first truncated level: the differential of multiplication by $p$
is zero.  Thus this is already a finite-stage frame of
$\omega_{\mathcal H_{\det}}$, and hence of $\mathcal N_s$.
Only this finite coefficient datum is needed; no entire infinite
marking is claimed to occur at one finite level.

The construction commutes with restrictions to successive label
boundaries and with their power identifications.  Alternatively,
pull back the determinant group from $\mathcal R_s$ first and
repeat the same finite first-level marking and splitting on that
stratum.  Both give the same line by base-change naturality.
Frobenius transports preserve $\nu_s$, whose values lie in
$\Fp^\times$.  Tensor powers and duals of a finite-stage frame
remain finite-stage frames.  This proves all the admissibility
assertions.
\end{proof}

\begin{corollary}[Coefficients with a norm-character filtration]
\label{cor:compact:norm-filtration}
Let $E$ be an admissible compact coefficient for a pair $(s,t)$,
and let $Q$ be a finite-dimensional smooth $P_s$-representation
with a finite $P_s$-stable filtration
\[
        0=Q_0\subset Q_1\subset\cdots\subset Q_l=Q,
        \qquad Q_j/Q_{j-1}\simeq k(\nu_s^{e_j}).
\]
Then $E\otimes_k Q$ and $E\otimes_k Q^\vee$ are admissible.
At every induction stage where compact finiteness at $(s,t)$ is
available, their total $P_s$-cohomology is finite-dimensional.
\end{corollary}

\begin{proof}
Each quotient of the filtration $E\otimes_k Q_j$ is
$E\otimes_k k(\nu_s^{e_j})$.  By
Lemma~\ref{lem:compact:norm-line}, tensoring with this constant
character is tensoring with a power of the admissible universal
determinant line.  On every required stratum take one common finite
frame and splitting level for the filtration of $E$ and the
finitely many determinant-line powers.  Refining the tensor
filtration then gives equivariantly trivial graded frames.
The original filtration is split as a module filtration, since it
is obtained by tensoring finite-dimensional $k$-vector spaces;
the refined filtration has the module splittings required in the
definition.  The same construction applies on every successive
boundary and after root transport.  Thus $E\otimes_k Q$ is
admissible.  Dualizing the finite filtration reverses its order
and replaces $e_j$ by $-e_j$, so it also proves the assertion for
$Q^\vee$.  Proposition~\ref{prop:compact-finiteness} gives the
cohomological conclusion.  All acting groups in this argument
are the full $P_s$.
\end{proof}

\subsection{Stable lattices and continuous duality}

\begin{lemma}[Lattices compatible with transfer]
\label{lem:compact:stable-lattices}
At a sufficiently deep finite joint frame level, the finite free
coefficient $B_U$ and the associated coefficients of any fixed finite
translation diagram admit $P_n$-stable free $C_1$-lattices.  Their fixed
pole strips have finite filtrations by boundary restrictions of
admissible coefficients.  For the stationary functions coefficient one
can choose such a lattice $L$ and an integer $b$ so that
\[
                    P=x^{-b}L
\]
is preserved by $S$, and every $x^{-c}L$ is carried into $P$ by a
sufficiently high iterate of $S$.  In the invariant-volume frames the
transpose on the dual lattice of top forms has Frobenius matrix entries
in $xC_1$.  Its finite basis, viewed as finite-cover functions by these
frames, can be taken integral over $C_1$.
\end{lemma}

\begin{proof}
Begin after the first \'{e}tale basis level and after the finite marking
level defining the volume and its character.  The tame connected
marking adjoins a $(p^t-1)$st root of the leading height coefficient,
with the line-frame convention understood.  Its coefficient module on
localization is a sum of powers of the invariant line.  Subsequent
normal finite levels can be taken to have $p$-group deck groups; their
regular representations over $k$ have finite filtrations by trivial
representations.  Associated-bundle descent gives filtrations of the
localized coefficient modules, split as module filtrations.  The same
construction applies to a finite translation diagram after one frame
level on which its extra finite actions are defined.  In the
multiplicative case one first uses its character decomposition.

Choose lifts of bases of the graded lines.  The action matrices in
these bases are triangular.  Their off-diagonal entries have one common
pole bound because the acting group is compact and the actions are
bounded-pole continuous.  Rescale successive lifted bases by powers of
$x$, starting at one end of the filtration.  At each step the finitely
many off-diagonal denominators are absorbed by the earlier rescalings.
The resulting split free lattice is stable.  Its successive quotients
are line lattices, and the same is true for its pole strips.
Factoring $x$ into label coordinates as in
Corollary~\ref{rings:boundary-labels} gives the asserted boundary
filtrations.

Choose a divisible stationary step whose Frobenius power fixes $k$,
and denote its coefficient Frobenius exponent by $Q>1$.  Finite
freeness of Frobenius and a finite root basis give a constant $c_0$ with
\[
 S(x^{-c}L)\subset
 x^{-\lceil c/Q\rceil-c_0}L.
\]
This is obtained by applying $S$ to the finitely many root-basis
vectors and clearing their denominators; inverse semilinearity then
divides the remaining pole order by $Q$.  If $b$ is larger than the
fixed point of this affine bound, $S(P)\subset P$, and iterating the
bound sends every fixed larger lattice into $P$.

In dual invariant-volume frames, Lemma~\ref{lem:cartier-transfer}
identifies the transpose with positive Frobenius.  Replacing the dual
lattice by $x^b$ times it increases its Frobenius matrix valuations by
$(Q-1)b$, up to the fixed denominators of its original matrix.  A
further increase of $b$ therefore makes every entry divisible by $x$.
Finally a finite basis of the finite \'{e}tale algebra over
$C_1[1/x]$ becomes integral over $C_1$ after multiplying by a large
power of $x$: clear the denominators in a monic equation for each basis
element.  Increase $b$ to achieve this simultaneously in the
volume/trace frames.  The preceding stability and positivity
inequalities continue to hold.
\end{proof}

\begin{lemma}[Compact and residue duality]\label{lem:compact:duality}
Let $L$ be a stable finite free lattice over a formal $r$-disc $C$, and
let $\mathfrak m_C$ be its closed-point ideal.  Put
$L^*=\Hom_C(L,C)$ and
$\omega_C=\bigwedge^r\widehat\Omega^1_{C/k}$, using continuous
differentials.
There are natural identifications
\begin{align*}
 L^\vee&\simeq H^r_{\mathfrak m_C}(L^*\otimes_C\omega_C),\\
 H^i(P_n,L)^\vee
 &\simeq H^{n^2-i}\bigl(P_n,L^\vee\bigr).
\end{align*}
Here $(-)^\vee=\Hom_{k,\mathrm{cts}}(-,k)$ is the discrete continuous
dual.  The identifications are compatible with finite coefficient
diagrams and with the transpose of the constant-cochain cup action.
In particular $H^i(P_n,L)$ is profinite and vanishes outside
$0\leq i\leq n^2$.
\end{lemma}

\begin{proof}
The units in the height-$n$ division algebra have analytic dimension
$n^2$.  They have no $p$-torsion in the present range: an element of
order $p$ would embed the degree-$p-1$ extension $\Qp(\zeta_p)$ in a
division algebra of degree $n$, which is impossible for $n<p-1$.
The compact analytic group results recalled in
\cite[Sections~1.1--1.2]{Venjakob} give a resolution of the trivial
compact $\Zp[[P_n]]$-module by finitely generated projectives, of
length $n^2$.  Reducing it modulo $p$ and extending scalars to $k$
gives the corresponding $k[[P_n]]$-resolution: the underlying
$\Zp$-modules are flat, so this base change is exact.  The residue-field
completed group ring is also Noetherian of finite global dimension
\cite[Proposition~3.3]{ArdakovWadsley}.  The dualizing row and its adjoint action are described in
\cite[Propositions~4.16 and~4.40, and Remark~4.23]{BGHS}.
Its orientation character is trivial here.  Indeed, after a splitting field the adjoint
representation is conjugation on a full matrix algebra, whose
determinant is one.  This is the orientable form of analytic
Poincar\'e duality; see \cite{Lazard,SerreGalois}.

The finite resolution computes continuous cohomology for compact
coefficients, by reduction to finite discrete quotients and inverse
limit.  The coefficient maps in those inverse systems are surjective
on resolution terms.  It agrees with the continuous bar model, and
also with that model on lattice unions using the specified pole
bounds.  For a compact lattice all differentials have closed image,
so cohomology is compact and dualization is exact.  Duality for the
finite projective resolution gives the second identification and the
cup-action compatibility.

For the first identification use formal coordinates $w_1,\ldots,w_r$.
The residue pairing sends a negative principal monomial times
$dw_1\wedge\cdots\wedge dw_r$ and a power series to the coefficient
of $w_1^{-1}\cdots w_r^{-1}$.  Every continuous functional on $C$
depends on a finite-dimensional power-series quotient, so it is a
finite sum of such monomial functionals.  Applying this pairing to a
free basis proves the assertion for $L$.  The residue transformation
law makes the pairing independent of coordinates and accounts for
the canonical line.  It therefore respects all the stated actions.
\end{proof}
